\chapter{État de l'art}

\section{Introduction}
Le présent chapitre vise à positionner le projet \textbf{Certif-fun} dans son contexte scientifique et technique, en s'appuyant sur un examen des approches et des travaux déjà réalisés dans le domaine. Il a pour rôle de situer notre démarche par rapport aux solutions existantes et d'aider à justifier la pertinence de notre proposition. Nous aborderons ici une recherche bibliographique sur les technologies de surveillance (\textit{Proctoring}), l'IA générative en éducation, ainsi qu'une comparaison des plateformes et architectures actuelles pour mettre en exergue les besoins encore non satisfaits.

\section{Travaux connexes et solutions existantes}
\subsection{Proctoring automatisé (IA de surveillance)}
De nombreux travaux démontrent l'efficacité de systèmes multi-modaux basés sur la vision par ordinateur. Par exemple, \textit{Singh et al.} proposent un proctoring IA qui utilise \textbf{YOLO} pour repérer les objets interdits (téléphones, livres) et \textbf{MediaPipe} pour suivre le visage et les yeux, en traitant en temps réel les flux vidéo des candidats. 

Le système détecte les téléphones et le nombre de personnes via \textbf{YOLOv8}, et peut automatiquement interrompre l'examen lorsqu'une fraude est détectée. D'autres solutions open-source, comme \textit{EduView}, intègrent des modèles TensorFlow et MediaPipe pour l'analyse des mouvements oculaires et utilisent Ultralytics YOLOv8 pour la détection faciale. Ces approches offrent modularité et efficacité sans nécessiter de ressources \textit{cloud}, en signalant en temps réel les activités suspectes.

\subsection{Outils de génération pédagogique par IA}
Plusieurs applications émergentes assistent les enseignants dans la création de quiz. Par exemple, des projets utilisant \textbf{FastAPI} et \textbf{LangChain} illustrent comment exploiter \textbf{Ollama} pour extraire le contenu de documents et générer automatiquement des questionnaires. 

Les recherches académiques corroborent ce potentiel : l'utilisation de grands modèles de langage (\textit{LLM}) peut produire un large éventail de questions alignées sur le cours. Toutefois, les essais montrent souvent une qualité inégale, soulignant que les quiz générés « couvrent un large spectre de sujets » mais peuvent manquer de profondeur dans les domaines de pointe sans un \textit{fine-tuning} approprié ou une supervision humaine.

\subsection{Plateformes existantes}
Certaines plateformes de certification intègrent déjà des fonctions avancées :
\begin{itemize}
    \item \textbf{HackerEarth :} Propose un environnement combinant gestion d'examens, verrouillage du navigateur et analyses anti-plagiat.
    \item \textbf{Moodle :} Le LMS \textit{open-source} s'enrichit désormais de plugins IA. La version 5.0 intègre un fournisseur Ollama permettant d'appeler des modèles LLM locaux pour générer ou résumer des contenus.
    \item \textbf{Solutions propriétaires :} Des outils comme \textit{Proctorio} ou \textit{Respondus} offrent du proctoring automatisé, mais restent fermés et coûteux.
\end{itemize}
Il manque toutefois une solution unifiée liant nativement proctoring et génération de contenu dans un seul système microservices.

\subsection{Architectures et technologies}
Les architectures microservices se démocratisent dans l'EdTech pour répondre aux besoins de scalabilité. Elles offrent un système modulaire où chaque service (gestion utilisateurs, cours, proctoring, visioconférence) est découplé. L'utilisation de \textbf{Spring Boot}, \textbf{FastAPI} et \textbf{React}, associée à \textbf{Docker}, permet une mise à jour isolée de chaque composant et facilite l'intégration d'algorithmes d'IA gourmands en ressources sans impacter le reste de la plateforme.

\section{Analyse comparative}
Cette partie compare les solutions citées selon divers critères techniques, économiques et fonctionnels.

\subsection{Synthèse comparative}
Le tableau \ref{table:comparaison_existant} résume les forces et faiblesses des approches actuelles par rapport à Certif-fun.

\begin{table}[H]
\centering
\caption{Comparaison des solutions existantes avec Certif-fun.}
\label{table:comparaison_existant}
\begin{tabular}{|p{3cm}|p{3cm}|p{3cm}|p{3cm}|}
\hline
\textbf{Critères} & \textbf{LMS Classiques (Moodle)} & \textbf{Proctoring Dédié (Proctorio)} & \textbf{Certif-fun} \\ \hline
\textbf{Surveillance IA} & Via plugins tiers & Très avancée & \textbf{Native (YOLOv8)} \\ \hline
\textbf{Génération IA} & Limitée / Externe & Aucune & \textbf{Native (Ollama)} \\ \hline
\textbf{Architecture} & Monolithique & Propriétaire & \textbf{Microservices} \\ \hline
\textbf{Coût} & Gratuit (Open source) & Élevé (Licence) & \textbf{Open Source / Local} \\ \hline
\textbf{Collaboration} & Basique & Absente & \textbf{Visioconférence intégrée} \\ \hline
\end{tabular}
\end{table}

\subsection{Analyse des aspects stratégiques}
\begin{itemize}
    \item \textbf{Aspects techniques :} L'association YOLOv8 + MediaPipe permet un traitement en temps réel sur du matériel standard, surpassant les versions antérieures en rapidité.
    \item \textbf{Aspects économiques :} L'usage d'Ollama (LLM local) réduit les coûts d'API cloud (OpenAI) et préserve la confidentialité des données (RGPD).
    \item \textbf{Aspects fonctionnels :} Certif-fun se distingue en unifiant sécurité et pédagogie, là où les autres solutions sont spécialisées dans un seul domaine.
\end{itemize}

\section{Conclusion}
Pour conclure ce chapitre, il est essentiel de faire une mise en perspective. L'analyse de l'état de l'art a permis de constater que si des solutions performantes existent pour la surveillance d'examens ou la gestion de cours, elles restent souvent incomplètes ou perfectibles car elles ne proposent pas d'intégration native et unifiée de l'intelligence artificielle générative et de la surveillance intelligente au sein d'une même architecture modulaire.

Cette étude de l'existant confirme et oriente nos choix stratégiques : l'adoption d'une architecture microservices pour la scalabilité, ainsi que l'utilisation de modèles locaux (\textit{Edge IA}) pour garantir à la fois performance et confidentialité des données. Les conclusions tirées ici serviront de base solide pour élaborer, dans le chapitre suivant, une architecture adaptée et une modélisation \textbf{UML} rigoureuse. Ce cadrage des approches préexistantes met en relief la valeur ajoutée de Certif-fun et pave la voie à la phase de \textbf{Conception} technique de notre solution.
